\documentclass[output=paper]{langscibook}
\ChapterDOI{10.5281/zenodo.22880284}
\title{\textit{A vaahoa kave goroto} – The launching of the turtle net}
\author{Enoch Horai Magum}
\abstract{}
\begin{document}
\maketitle\label{chap:28}
\begin{xltabular}{\textwidth}{lXX}
	(\ref{ex:text21-1})&\textit{A vaahoa kave goroto na tei me nana bona maa taatate nae.}&The launching of the turtle net is associated with certain customs.\\
	(\ref{ex:text21-2})&\textit{O taatate bona o vasisio na vavasisio kurusu ni raara.}&These respected customs are very much respected by us.\\
	(\ref{ex:text21-3})&\textit{A otei vai a too aito toro viana ta buaku ge ta dorana ta otei, eove repaa mee ri bari rasuu tea aito.}&A man, an expert on herbs, selects two or four men, and then he takes them into the jungle to collect herbs.\\
	(\ref{ex:text21-4})&\textit{Beori aito vakavara maa, eori repaa ma hee ni bona aito tea peha moon to pasi ato ge kikira ni nom bona.}&When they have collected the herbs, they give them to a woman who will hold them or take care of them.\\
	(\ref{ex:text21-5})&\textit{A aito vai toro paku ara, beara vaasubunubu ni rara a vaahoa na kave goroto.}&We must collect these herbs, when we plan the launching of the turtle net for the next day.\\
	(\ref{ex:text21-6})&\textit{A moon vai to kahi ato ni nana bona aito toro vaavavina nana bona too vaahito na kave goroto.}&The woman who will hold the herbs must be the sister of the owner of the turtle net.\\
	(\ref{ex:text21-7})&\textit{Teo bon vaahoa na kave, eara repaa isuvu to buo koopu ruene,}&On the day of the net launch, we fill two bamboo water containers,\\
	(\ref{ex:text21-8})&\textit{eara repaa vaatei a aito teebona.}&and then we put the herbs into them.\\
	(\ref{ex:text21-9})&\textit{Beara tau tahoa rakaha, eara repaa vaobete o buo ruene komano sinivi to vaahae batana bona kave.}&Just when we are about to set off, we put the two water containers into the canoe which has the net on board.\\
	(\ref{ex:text21-10})&\textit{Beara nao komana tahii, eara repaa sisi mana a kave teebona.}&When we put to sea, we sprinkle the net with them.\\
	(\ref{ex:text21-11})&\textit{O hum busubusu ge veo goroto na vanata komana ni raara ei saka tea irihi haana,}&The place where the net is thrown or where the turtle netting is done, we know, this is not at the edge of the reef.\\
	(\ref{ex:text21-12})&\textit{Ahiki, amaa hum veo goroto toro kakaamuru to hiki nae o vasu.}&No, the places for turtle netting must be the white sand without stones.\\
	(\ref{ex:text21-13})&\textit{Beara busu gunaha nia kave en teo hum vasuvasu, a kave toro hana.}&If we threw the net at a place full of stones, the net would get caught up there.\\
	(\ref{ex:text21-14})&\textit{Mea maa hum veo goroto teara a ta Teapu ei Taenaruene, Teopasuna, Teanin.}&The turtle netting places of us Teop people are \textit{Taenaruene}, \textit{Teopasuna}, \textit{Teanin}.\\
	(\ref{ex:text21-15})&\textit{Amaa hum bona kou na hiki nana to vasu, o kakaamuru koa.}&These places do not have any stones, just white sand.\\
\end{xltabular}
\end{document}